% ============================================================
\chapter{Continuous Random Variables}
\label{ch:continuous}
% ============================================================

With discrete variables, we counted. With continuous variables, we measure. The transition is more than a change of vocabulary: it requires replacing sums with integrals, point masses with densities, and coming to terms with the idea that in the continuous case, the probability of any \emph{exact} value is always zero. It was Abraham de Moivre who, around 1733, opened this path by approximating the binomial distribution with the bell curve---what we now call the normal distribution, and which Gauss would immortalise in his astronomical work.

\begin{intuition}
A continuous random variable takes values in an interval (or all of $\R$) and its
distribution is described by a \textbf{probability density function}. This chapter
develops the tools specific to the continuous case: densities, cumulative
distribution functions, and quantiles.
\end{intuition}

% ------------------------------------------------------------
\section{Probability Density Functions}
% ------------------------------------------------------------

\begin{definition}[Absolutely continuous random variable]
A random variable $X$ is \textbf{continuous} (or absolutely continuous) if there
exists an integrable function $f_X : \R \to [0, +\infty)$ such that for all $a \leq b$:
\[
    \PP(a \leq X \leq b) = \int_a^b f_X(x)\, dx.
\]
The function $f_X$ is called the \textbf{probability density function} (pdf) of $X$.
\end{definition}

\begin{proposition}[Properties of the pdf]
\begin{enumerate}[label=(\roman*)]
    \item $f_X(x) \geq 0$ for all $x \in \R$.
    \item $\int_{-\infty}^{+\infty} f_X(x)\, dx = 1$.
    \item $\PP(X = a) = 0$ for every $a \in \R$.
    \item $\PP(a \leq X \leq b) = \PP(a < X < b) = \PP(a \leq X < b) = \PP(a < X \leq b)$.
\end{enumerate}
\end{proposition}

\begin{warning}[title=\textbf{A density is not a probability}]
It is possible for $f_X(x) > 1$ at some values of $x$. The density is a
``concentration of probability'', not a probability itself. Only its integral
over an interval gives a probability.
\end{warning}

% ------------------------------------------------------------
\section{Cumulative Distribution Function}
% ------------------------------------------------------------

\begin{definition}[CDF]
\[
    F_X(x) = \PP(X \leq x) = \int_{-\infty}^{x} f_X(t)\, dt, \quad x \in \R.
\]
\end{definition}

\begin{theorem}[Density--CDF relationship]
If $F_X$ is differentiable at $x$, then $f_X(x) = F_X'(x)$. More generally,
$f_X$ equals the derivative of $F_X$ almost everywhere.
\end{theorem}

\begin{proposition}[Properties of $F_X$]
\begin{enumerate}[label=(\roman*)]
    \item $F_X$ is non-decreasing and continuous (in the absolutely continuous case).
    \item $\lim_{x \to -\infty} F_X(x) = 0$, $\lim_{x \to +\infty} F_X(x) = 1$.
    \item $\PP(a < X \leq b) = F_X(b) - F_X(a)$.
\end{enumerate}
\end{proposition}

% ------------------------------------------------------------
\section{Classical Continuous Distributions}
% ------------------------------------------------------------

\subsection{Continuous uniform distribution}

\begin{definition}[Uniform on $\lbrack a,b\rbrack$]
$X \sim \mathcal{U}([a,b])$ with $a < b$:
\[
    f_X(x) = \frac{1}{b-a}\,\mathbf{1}_{[a,b]}(x), \quad
    F_X(x) = \begin{cases} 0 & \text{if } x < a, \\ \frac{x-a}{b-a} & \text{if } a \leq x \leq b, \\ 1 & \text{if } x > b. \end{cases}
\]
$\E[X] = (a+b)/2$, $\Var(X) = (b-a)^2/12$.
\end{definition}

\subsection{Exponential distribution}

\begin{definition}[Exponential]
$X \sim \mathrm{Exp}(\lambda)$ with $\lambda > 0$:
\[
    f_X(x) = \lambda e^{-\lambda x}\,\mathbf{1}_{[0,+\infty)}(x), \quad
    F_X(x) = (1 - e^{-\lambda x})\,\mathbf{1}_{[0,+\infty)}(x).
\]
$\E[X] = 1/\lambda$, $\Var(X) = 1/\lambda^2$.
\end{definition}

\begin{theorem}[Memorylessness]
The exponential distribution is the only continuous distribution on $[0,+\infty)$
satisfying:
\[
    \PP(X > s + t \mid X > s) = \PP(X > t) \quad \forall\, s, t \geq 0.
\]
\end{theorem}

\begin{proof}
Let $g(t) = \PP(X > t) = e^{-\lambda t}$. Then:
\[
    \PP(X > s+t \mid X > s) = \frac{e^{-\lambda(s+t)}}{e^{-\lambda s}}
    = e^{-\lambda t} = g(t).
\]
Conversely, if $g(s+t) = g(s)g(t)$ with $g$ continuous, decreasing, and $g(0)=1$,
then $g(t) = e^{-\lambda t}$ for some $\lambda > 0$, which characterises the
exponential distribution.
\end{proof}

\begin{intuition}[title=\textbf{Exponential and Poisson}]
If arrivals follow a Poisson process of rate $\lambda$, then the time between
consecutive arrivals follows an $\mathrm{Exp}(\lambda)$ distribution.
\end{intuition}

\subsection{Normal (Gaussian) distribution}

\begin{definition}[Normal distribution]
$X \sim \mathcal{N}(\mu, \sigma^2)$ with $\mu \in \R$ and $\sigma > 0$:
\[
    f_X(x) = \frac{1}{\sigma\sqrt{2\pi}}
    \exp\!\left(-\frac{(x-\mu)^2}{2\sigma^2}\right), \quad x \in \R.
\]
$\E[X] = \mu$, $\Var(X) = \sigma^2$.
\end{definition}

\begin{definition}[Standard normal]
If $Z \sim \mathcal{N}(0,1)$, we write $\Phi(x) = \PP(Z \leq x)$ for its CDF.
The standardisation relationship is:
\[
    \text{if } X \sim \mathcal{N}(\mu, \sigma^2), \quad
    \PP(X \leq x) = \Phi\!\left(\frac{x - \mu}{\sigma}\right).
\]
\end{definition}

\begin{keyformulas}[title=\textbf{Key values of $\Phi$}]
\begin{center}
\renewcommand{\arraystretch}{1.3}
\begin{tabular}{cccccc}
    \toprule
    $z$ & 1.645 & 1.96 & 2.326 & 2.576 & 3.291 \\
    \midrule
    $\Phi(z)$ & 0.95 & 0.975 & 0.99 & 0.995 & 0.9995 \\
    \bottomrule
\end{tabular}
\end{center}
By symmetry: $\Phi(-z) = 1 - \Phi(z)$.
\end{keyformulas}

\begin{proposition}[The $k\sigma$ rule]
If $X \sim \mathcal{N}(\mu, \sigma^2)$:
\begin{itemize}
    \item $\PP(\abs{X - \mu} \leq \sigma) \approx 0.6827$ \quad ($1\sigma$ rule),
    \item $\PP(\abs{X - \mu} \leq 2\sigma) \approx 0.9545$ \quad ($2\sigma$ rule),
    \item $\PP(\abs{X - \mu} \leq 3\sigma) \approx 0.9973$ \quad ($3\sigma$ rule).
\end{itemize}
\end{proposition}

\subsection{Gamma distribution}

\begin{definition}[Gamma]
$X \sim \Gamma(\alpha, \lambda)$ with $\alpha, \lambda > 0$:
\[
    f_X(x) = \frac{\lambda^\alpha}{\Gamma(\alpha)}\,x^{\alpha-1}\,e^{-\lambda x}\,
    \mathbf{1}_{(0,+\infty)}(x),
\]
where $\Gamma(\alpha) = \int_0^{\infty} t^{\alpha-1}e^{-t}\,dt$. We have
$\E[X] = \alpha/\lambda$ and $\Var(X) = \alpha/\lambda^2$.
\end{definition}

\begin{remark}
$\Gamma(1, \lambda) = \mathrm{Exp}(\lambda)$ and $\Gamma(n/2, 1/2) = \chi^2(n)$
(the chi-squared distribution with $n$ degrees of freedom).
\end{remark}

\subsection{Beta distribution}

\begin{definition}[Beta]
$X \sim \mathrm{Beta}(\alpha, \beta)$ with $\alpha, \beta > 0$:
\[
    f_X(x) = \frac{x^{\alpha-1}(1-x)^{\beta-1}}{B(\alpha,\beta)}\,
    \mathbf{1}_{(0,1)}(x),
\]
where $B(\alpha,\beta) = \Gamma(\alpha)\Gamma(\beta)/\Gamma(\alpha+\beta)$.
$\E[X] = \alpha/(\alpha+\beta)$.
\end{definition}

% ------------------------------------------------------------
\section{Quantiles and Median}
% ------------------------------------------------------------

\begin{definition}[Quantile]
The \textbf{quantile of order $p$} ($0 < p < 1$) of $X$ is:
\[
    q_p = \inf\{x \in \R : F_X(x) \geq p\} = F_X^{-1}(p).
\]
\end{definition}

\begin{definition}[Median]
The \textbf{median} of $X$ is the quantile of order $1/2$:
$\mathrm{Med}(X) = q_{1/2}$.
\end{definition}

\begin{example}
If $X \sim \mathrm{Exp}(\lambda)$, then $F_X(x) = 1 - e^{-\lambda x}$ for $x \geq 0$.
The quantile of order $p$ is:
\[
    q_p = -\frac{\ln(1-p)}{\lambda}.
\]
In particular, $\mathrm{Med}(X) = \ln 2/\lambda$.
\end{example}

% ------------------------------------------------------------
\section{Simulation by Inversion}
% ------------------------------------------------------------

\begin{theorem}[Inverse transform method]
If $U \sim \mathcal{U}([0,1])$ and $F$ is a continuous, strictly increasing CDF,
then $X = F^{-1}(U)$ has CDF $F$.
\end{theorem}

\begin{proof}
$\PP(X \leq x) = \PP(F^{-1}(U) \leq x) = \PP(U \leq F(x)) = F(x)$.
\end{proof}

\begin{example}
To simulate $X \sim \mathrm{Exp}(\lambda)$, generate $U \sim \mathcal{U}([0,1])$
and set $X = -\ln(1-U)/\lambda = -\ln(U)/\lambda$ (since $1-U$ has the same
distribution as $U$).
\end{example}

% ------------------------------------------------------------
\section{Exercises}
% ------------------------------------------------------------

\begin{exercise}
Let $X$ have density $f(x) = cx^2\,\mathbf{1}_{[0,1]}(x)$. Find $c$, $\E[X]$,
$\Var(X)$, and $F_X$.
\end{exercise}

\begin{exercise}
Let $X \sim \mathcal{N}(100, 25)$. Compute:
\begin{enumerate}[label=(\alph*)]
    \item $\PP(X > 105)$;
    \item $\PP(90 < X < 110)$;
    \item the quantile $q_{0.9}$.
\end{enumerate}
\end{exercise}

\begin{exercise}
Show that if $X \sim \mathcal{N}(\mu, \sigma^2)$, then $Z = (X - \mu)/\sigma
\sim \mathcal{N}(0,1)$.
\end{exercise}

\begin{exercise}
Let $X \sim \mathrm{Exp}(\lambda)$. Find the density and expectation of $Y = X^2$.
\end{exercise}

\begin{exercise}
Show that the Gaussian integral equals
$\int_{-\infty}^{+\infty} e^{-x^2/2}\,dx = \sqrt{2\pi}$.
\textit{Hint:} compute $I^2$ using polar coordinates.
\end{exercise}

\begin{exercise}
Let $X \sim \Gamma(\alpha, \lambda)$. Show that $Y = 2\lambda X \sim \chi^2(2\alpha)$
when $\alpha \in \N^*/2$.
\end{exercise}

\begin{exercise}[Simulation]
Describe a method to simulate a random variable with density
$f(x) = 2x\,\mathbf{1}_{[0,1]}(x)$ using the inverse transform method.
\end{exercise}

\begin{exercise}
Let $X$ have density $f(x) = \frac{2}{\pi(1+x^2)}\,\mathbf{1}_{(0,+\infty)}(x)$.
Verify that $f$ is a valid density and compute $F_X$. Does $\E[X]$ exist?
\end{exercise}
